% ID: SEL-D0068-OB-E02
% ESTADO: cotejada. Segunda pasada visual y borrador compilado revisados.
% Fuente: Selectividad/sel_2009_matematicas_aplicadas/mat apli CC SS02.pdf
% SHA-256: 24ee6d89c206838cb0cf4f32c4e54391e7ca81f77c4dd0be651ad7218f37beb6
% Página: 2; opción B; ejercicio original 2.

Sea la función $f(x)=ax^3+bx^2+x$.
\begin{enumerate}[label=\alph*)]
\item \textbf{(1.5 puntos)} Determine el valor de los parámetros $a$ y $b$ sabiendo que la función $f$ tiene un máximo en $x=1$ y que $f(1)=2$.
\item \textbf{(1.5 puntos)} Para $a=b=1$, halle la ecuación de la recta tangente a la gráfica de $f$ en el punto de abscisa $x=0$.
\end{enumerate}
